\begin{definition}
В этом билете граф простой, \(n\ge3\),
\(d_1\le\cdots\le d_n\) — упорядоченные степени.
Гамильтоновость означает наличие простого цикла через все вершины.
\end{definition}
\begin{theorem}[Оре]
Если для любых двух несмежных вершин \(u,v\)
\(d(u)+d(v)\ge n\), граф гамильтонов.
\end{theorem}
\begin{proof}
Выведем утверждение из критерия Хватала.


Предположим, что критерий Хватала нарушен при \(1\le i<n/2\):
\(d_i\le i\), \(d_{n-i}\le n-i-1\). Пусть \(T\) — первые \(n-i\)
вершин упорядоченного списка, \(S\subset T\) — первые \(i\).
Для \(x\in S\) и \(y\in T\setminus\{x\}\) сумма степеней не превосходит
\(i+(n-i-1)=n-1\). По условию Оре \(x,y\) обязательно смежны.
Отсюда \(i\ge d(x)\ge n-i-1\ge i\), где последнее неравенство
следует из \(n>2i\). Все неравенства должны быть равенствами:
\(n=2i+1\), \(d(x)=i\), и у \(x\) нет соседей вне \(T\).
В \(T\setminus S\) осталась ровно одна вершина \(z\).
Она смежна со всеми \(i\) вершинами \(S\) и имеет степень не больше \(i\),
так что тоже не имеет соседей вне \(T\).
Получили компоненту \(T=K_{i+1}\); дополнение размера \(i\) непусто
и его вершины имеют степень не больше \(i-1\).
Выбрав \(x\in T\), \(w\notin T\), получаем несмежную пару с суммой
степеней не больше \(2i-1<n\), противоречие Оре.
Следовательно, условие Хватала выполнено и граф гамильтонов.
\end{proof}
Есть и более короткое доказательство через замыкание из билета 10:
любую недостающую пару можно дополнить ребром, условие Оре при этом
сохраняется, в итоге получаем \(K_n\). Лемма о добавлении ребра
сохраняет гамильтоновость в обе стороны.
\begin{theorem}[Дирак]
Если \(\delta(G)\ge n/2\), то граф гамильтонов.
\end{theorem}
\begin{proof}
Для любых несмежных \(u,v\) имеем \(d(u)+d(v)\ge2\delta(G)\ge n\).
Применима теорема Оре. Непосредственно условие Хватала тоже выполнено:
при \(i<n/2\) имеем \(d_i\ge\delta(G)>i\), поэтому посылка импликации ложна.
\end{proof}
\begin{cor}[Версии для пути]
Если \(n\ge2\) и сумма степеней любой несмежной пары не меньше \(n-1\),
есть гамильтонов путь. Достаточно также условия \(\delta(G)\ge(n-1)/2\).
\end{cor}
\begin{proof}
Добавим универсальную вершину \(w\). В новом графе \(n+1\) вершина,
а степени старой несмежной пары в сумме не меньше \((n-1)+2=n+1\).
У \(w\) несмежных пар нет, поэтому применима теорема Оре.
Удаление \(w\) из гамильтонова цикла оставляет гамильтонов путь.
Условие на минимальную степень влечёт условие на сумму.
\end{proof}
\begin{example}[Оценку Дирака нельзя понизить на единицу]
При \(n=2r\), \(r\ge2\), граф \(K_{r-1,r+1}\) имеет минимальную степень
\(r-1\), но не имеет гамильтонова цикла: в двудольном цикле числа
вершин обеих долей равны. Условия теорем являются достаточными;
их нарушение само по себе не доказывает негамильтоновость.
\end{example}
